% Brief topic description for ONR Program Officer discussion (FY27 ONR YIP, N0001426SF004).
% Compile with XeLaTeX from the repository root:
%   latexmk -xelatex -cd -interaction=nonstopmode ONR/claude_src/claude_behzad_kamgar_parsi.tex
\documentclass[11pt]{article}
\usepackage[letterpaper,top=0.55in,bottom=0.5in,left=0.72in,right=0.72in]{geometry}
\usepackage{fontspec}
\IfFileExists{ben/fonts/EBGaramond-Regular.ttf}{\def\cnfont{ben/fonts/}}{%
  \IfFileExists{../../ben/fonts/EBGaramond-Regular.ttf}{\def\cnfont{../../ben/fonts/}}{\def\cnfont{../ben/fonts/}}}
\setmainfont{EBGaramond-Regular.ttf}[Path=\cnfont,BoldFont=EBGaramond-Bold.ttf,
  ItalicFont=EBGaramond-Italic.ttf,BoldItalicFont=EBGaramond-BoldItalic.ttf,Numbers=Lining]
\usepackage{microtype}
\usepackage{xcolor}
\definecolor{navy}{HTML}{1F3A5F}
\usepackage[colorlinks=true,urlcolor=navy,linkcolor=navy]{hyperref}
\usepackage{enumitem}
\pagestyle{empty}
\setlength{\parindent}{0pt}
\setlength{\parskip}{3.5pt}
\emergencystretch=2em
\newcommand{\head}[1]{{\color{navy}\textbf{#1}}\enspace}
\setlist[itemize]{leftmargin=1.1em,itemsep=2pt,topsep=1.5pt,parsep=0pt}
\hypersetup{pdftitle={Self-Assessing AI Agents},pdfauthor={Xinya Du}}

% ---- Edit the Program Officer block here --------------------------------------
\newcommand{\POline}{Dr.\ Behzad Kamgar-Parsi, ONR Code 311, Artificial Intelligence (Machine Learning, Reasoning and Intelligence)}
\newcommand{\FocusArea}{Artificial Intelligence and Autonomy --- Machine Learning; Intelligent and Collaborative Agents}
% -------------------------------------------------------------------------------

\begin{document}
\begin{center}
{\fontsize{15.5}{18.5}\selectfont\bfseries\color{navy} Self-Assessing AI Agents: Learning Where Their Own Reasoning Fails\\[1pt] and What to Do Next\par}
\vspace{3pt}
{\small\itshape Brief topic description for Program Officer discussion \,·\, FY27 ONR Young Investigator Program (N0001426SF004)\par}
\end{center}
\vspace{-2pt}
{\small
\textbf{PI:} Xinya Du, Assistant Professor of Computer Science, The University of Texas at Dallas \,·\, \href{mailto:xinya.du@utdallas.edu}{xinya.du@utdallas.edu} \,·\, \href{https://xinyadu.github.io}{xinyadu.github.io}\\
\textbf{Program Officer:} \POline\\
\textbf{Proposed S\&T Focus Area:} \FocusArea\par}
\vspace{-3pt}
{\color{navy}\rule{\linewidth}{0.6pt}}

\head{Problem.} Agents built on large language models (LLMs) now search documents, query databases, read images, and call tools over many steps with little human supervision. They are still trained mostly on outcomes: a trajectory is rewarded if its final answer is right. Outcome feedback cannot say which step failed or why, and generic self-critique often blames the wrong step. An agent checking whether a resupply shipment covers a ship's request may query the supply database correctly but read the quantity from the wrong line of a photographed packing list. An untyped step score says only that step 3 looks doubtful. The useful diagnosis, \emph{visual entity error at step 3}, also says what to do: re-examine the photo, not re-query the database.

\head{Central hypothesis.} Agents can learn \emph{self-assessment}: calibrated judgments of which of their own steps are unsupported and what kind of error each is. We hypothesize that (i) verifiers trained on typed, evidence-linked feedback, synthesized by injecting errors of known type into correct trajectories, will localize agent errors more accurately than outcome rewards or untyped step scores; and (ii) agents that choose their next action from the diagnosed error type will succeed more often per tool call than agents that retry or self-critique generically.

\head{Research thrusts.}
\begin{itemize}
\item \textbf{Typed process verifiers for agent trajectories.} Represent each step as a claim, the evidence it relies on (passage, database record, image region, or tool output), and its preconditions. Extend our FG-PRM error-injection method from mathematical reasoning to agent trajectories: insert fabricated or misread evidence, invalid inferences, violated preconditions, wrong tool arguments, and visual entity, attribute, and relation errors, then train verifiers to detect, type, and localize them. Measure which synthetic error distributions transfer to the agent's natural errors.
\item \textbf{From diagnosis to action.} Use the diagnosis to choose the next step (re-read a source, re-examine an image region, re-call a tool with a corrected argument, backtrack, or abstain and ask an operator) while penalizing unnecessary retries. Calibrate the verifier's own confidence so the agent knows when its self-assessment is unreliable, and characterize when self-assessment helps planning and when it induces over-caution.
\item \textbf{Learning from self-assessment.} Use typed step-level feedback as dense rewards for reinforcement learning (RL), extending our hierarchical process rewards for agentic search (HiPRAG), and as labels for self-training on unlabeled trajectories. Compare training-time and inference-time use of the same verifier at equal computation.
\end{itemize}

\head{What is new.} Process reward models are trained mostly on mathematics with untyped labels, and agents mostly on outcome rewards. Here, self-assessment is typed, tied to evidence in text, databases, tools, and images, and used to act and to learn.

\head{Evaluation.} Multi-hop and event-centric question answering (MuSiQue; our MEQA benchmark), multimodal science questions (ScienceQA), and controlled tool-use tasks with injected argument errors and held-out tools. Measures: error detection, typing, and localization; calibration of self-assessed confidence; task success per tool call; and repair success. Baselines: outcome-reward RL, untyped process reward models, and LLM self-critique.

\head{Naval relevance.} Navy logistics and planning assistants will combine documents, databases, imagery, and calculators with limited human oversight, where an undetected wrong step is costly. An agent that can report \emph{step 3 read the quantity from the wrong line} can request a targeted check or hand the decision to an operator instead of failing silently. The topic addresses the program's interests in self-assessment, reasoning with uncertain information, and integrating knowledge, learning, and planning.

\head{PI foundation.} \href{https://doi.org/10.18653/v1/2025.findings-emnlp.228}{FG-PRM}, process reward models trained on synthetic hallucinations (EMNLP Findings 2025); \href{https://arxiv.org/abs/2510.07794}{HiPRAG}, hierarchical process rewards for efficient agentic retrieval (ICLR 2026); \href{https://doi.org/10.18653/v1/2024.findings-emnlp.290}{FaithScore} and \href{https://arxiv.org/abs/2404.05046}{FGAIF}, fine-grained hallucination evaluation and AI feedback for vision-language models (EMNLP Findings 2024; TMLR 2025); \href{https://proceedings.neurips.cc/paper_files/paper/2024/hash/e560a0b22e4432003d0dba63ff8dc457-Abstract.html}{MEQA} (NeurIPS 2024). \href{https://www.nsf.gov/awardsearch/showAward?AWD_ID=2340435}{NSF CAREER} award (2024); \href{https://www.amazon.science/research-awards/recipients/xinya-du}{Amazon Research Award} (2023) for process-guided fine-tuning. As a postdoc, the PI led system integration for the RESIN team in DARPA KAIROS (\href{https://doi.org/10.18653/v1/2022.naacl-demo.7}{\mbox{RESIN-11}}, NAACL 2022 Demo, first author).

\end{document}
